% This file is part of JTS-MM (https://github.com/paulotto/jts-mm).
% SPDX-FileCopyrightText: 2026 Paul-Otto Müller
% SPDX-License-Identifier: CC-BY-4.0

\section{Reproducing and editing the figures}\label{sec:reproduce}
Everything in this document is generated from \texttt{../../scan/teeth\_02.stl} by the scripts in
\texttt{scripts/}. Install the Python packages in \texttt{scripts/requirements.txt} and run
\texttt{scripts/run\_all.sh}; it rebuilds all figures, tables, STLs and this PDF (about
\SI{3}{min}; the 3D renders use VTK off-screen and need no display).

\begin{table}[H]\centering\small
\begin{tabularx}{\textwidth}{l X l}
\toprule
script & what it does & output\\
\midrule
\texttt{s00\_export\_teeth\_freecad.py} & optional, run with \texttt{freecadcmd}: exports \texttt{Teeth\_02} with its placement & \texttt{../../scan/}\\
\texttt{s01\_arch\_frame.py} & arch centreline, local frame, occlusal profile & \texttt{build/arch.npz}\\
\texttt{s02\_voxelize.py} & filled \SI{0.2}{mm} voxel model, distance field, arch coordinates & \texttt{build/field.npz}\\
\texttt{s03\_survey.py} & height of contour and undercut per tooth & \texttt{build/survey.json}, \cref{tab:survey}\\
\texttt{s04\_concepts.py} & the five concepts as voxel masks + wire paths, built from \texttt{params.py} & \texttt{stl/}, \texttt{build/concepts.pkl}\\
\texttt{s05\_render3d.py} & 3D renders, cameras in \texttt{CAMS} & \texttt{figures/png/render\_*}\\
\texttt{s06\_fig\_mechanics.py} & \cref{fig:mech,fig:undercut} & \texttt{mech\_*}\\
\texttt{s07\_fig\_sections.py} & first-molar sections of every concept & \texttt{section\_*}\\
\texttt{s08\_springs.py} & spring mathematics, \cref{fig:springs,fig:loopchart}, \cref{tab:wire,tab:loop,tab:tpu,tab:energy,tab:arm} & \texttt{springs\_*}, \texttt{generated/}\\
\texttt{s09\_drawings.py} & dimensioned drawings, reference dimensions, parameter tables & \texttt{drw\_*}, \texttt{generated/}\\
\texttt{s10\_export\_cad.py} & reference curves, stations, tooth sections and the FreeCAD macro & \texttt{cad/}\\
\texttt{params.py} & \textbf{every part dimension} (read by s04, s08, s09) & --\\
\bottomrule
\end{tabularx}
\end{table}

\paragraph{Editing figures.} Every 2D figure is written three times: \texttt{figures/svg/} (text is
kept as text, so labels can be edited in Inkscape or Illustrator), \texttt{figures/pdf/} (used by
this document) and \texttt{figures/png/} (preview). If you edit an SVG by hand, export it to PDF
under the same name (Inkscape: \emph{Save a Copy} $\rightarrow$ PDF) and LaTeX picks it up; but a
re-run of the scripts overwrites manual edits, so prefer changing the script where possible. The 3D
renders are PNG only: a vector export of \num{e5} triangles is neither small nor editable. To
change what they show, edit \texttt{params.py} (dimensions) or the masks in \texttt{s04\_concepts.py} or the camera in
\texttt{s05\_render3d.py}.

\paragraph{Coordinates.} $x$ mesio-distal, $y$ arch depth, $z$ tooth axis with the cusp tips at the
smallest $z$ (the frame of \texttt{Teeth\_02} in the FreeCAD files of the parts). The
scripts add $h$ (height below the lowest cusp tip), $s$ (arc length along the arch, 0 at the
midline) and $n$ (offset from the arch centreline, positive buccal). Every mask in
\texttt{s04\_concepts.py} is written in $s$, $n$, $h$ and the distance to the teeth.
